\documentclass[11pt]{article}
\usepackage[utf8]{inputenc}
\usepackage{amsmath,amssymb}
\usepackage{hyperref}
\title{Efficient Long Read Sequencing Assembly: A Resource-Aware Approach}
\author{Anonymous}
\date{}
\begin{document}
\maketitle

\begin{abstract}
This paper studies Long Read Sequencing Assembly. Grounded in a real, citable literature \cite{ref1,ref2,ref3,ref4,ref5,ref6,ref7,ref8},
we propose a focused method and report \textbf{illustrative} experiments.
All references below are real and verifiable (OpenAlex/arXiv); the reported
numbers are simulated to illustrate intended behaviour.
\end{abstract}

\section{Introduction}
Long Read Sequencing Assembly is an active research area. This work targets a concrete gap identified from
the recent literature and contributes a method together with an evaluation protocol.

\section{Related Work (real literature)}
The following works, retrieved from public bibliographic APIs, frame our study:
\begin{itemize}
\item Li (2018) — "Minimap2: pairwise alignment for nucleotide sequences", Bioinformatics (cited 17224).
\item Magoč et al. (2011) — "FLASH: fast length adjustment of short reads to improve genome assemblies", Bioinformatics (cited 15908).
\item Wick et al. (2017) — "Unicycler: Resolving bacterial genome assemblies from short and long sequencing reads", PLoS Computational Biology (cited 8935).
\item Koren et al. (2017) — "Canu: scalable and accurate long-read assembly via adaptive <i>k</i> -mer weighting and repeat separation", Genome Research (cited 8269).
\item Kolmogorov et al. (2019) — "Assembly of long, error-prone reads using repeat graphs", Nature Biotechnology (cited 6480).
\item Jonathan et al. (2020) — "Denoising Diffusion Probabilistic Models", arXiv (Cornell University) (cited 5653).
\end{itemize}

\section{Method}
We reduce the dominant computational cost via adaptive resource allocation, activating only the components needed per input. We instantiate this idea for Long Read Sequencing Assembly and analyse its cost/quality trade-off.

\section{Experiments (Illustrative)}
\begin{center}
\begin{tabular}{lcc}
System & Primary Metric & Efficiency \\ \hline
Strong baseline & 0.812 & 1.00$\times$ \\
Ablation & 0.828 & 0.92$\times$ \\
\textbf{Ours} & \textbf{0.851} & \textbf{0.71$\times$}
\end{tabular}
\end{center}
\emph{Metrics simulated to illustrate the intended effect; not measured.}

\section{Conclusion}
We connected a real literature to a concrete method for Long Read Sequencing Assembly. Future work will
validate the illustrative results with real experiments.

\begin{thebibliography}{9}
\bibitem{ref1} Heng Li. Minimap2: pairwise alignment for nucleotide sequences. \emph{Bioinformatics}, 2018. DOI:10.1093/bioinformatics/bty191
\bibitem{ref2} Tanja Magoč, Steven L. Salzberg. FLASH: fast length adjustment of short reads to improve genome assemblies. \emph{Bioinformatics}, 2011. DOI:10.1093/bioinformatics/btr507
\bibitem{ref3} Ryan R. Wick, Louise M. Judd, Claire L. Gorrie et al.. Unicycler: Resolving bacterial genome assemblies from short and long sequencing reads. \emph{PLoS Computational Biology}, 2017. DOI:10.1371/journal.pcbi.1005595
\bibitem{ref4} Sergey Koren, Brian P. Walenz, Konstantin Berlin et al.. Canu: scalable and accurate long-read assembly via adaptive <i>k</i> -mer weighting and repeat separation. \emph{Genome Research}, 2017. DOI:10.1101/gr.215087.116
\bibitem{ref5} Mikhail Kolmogorov, Jeffrey Yuan, Yu Lin et al.. Assembly of long, error-prone reads using repeat graphs. \emph{Nature Biotechnology}, 2019. DOI:10.1038/s41587-019-0072-8
\bibitem{ref6} Ho, Jonathan, Ajay N. Jain, Pieter Abbeel. Denoising Diffusion Probabilistic Models. \emph{arXiv (Cornell University)}, 2020. DOI:10.48550/arxiv.2006.11239
\bibitem{ref7} Ruibang Luo, Binghang Liu, Yinlong Xie et al.. SOAPdenovo2: an empirically improved memory-efficient short-read <i>de novo</i> assembler. \emph{GigaScience}, 2012. DOI:10.1186/2047-217x-1-18
\bibitem{ref8} Heng Li. Aligning sequence reads, clone sequences and assembly contigs with BWA-MEM. \emph{DROPS (Schloss Dagstuhl – Leibniz Center for Informatics)}, 2013. DOI:10.48550/arxiv.1303.3997
\end{thebibliography}
\end{document}
